% PAGES: 170-173
% Continues after the c2 chunk (PDF pages 163-169 / folios 147-153). c2's
% last file (sec05_c2_4.tex) ends cleanly on page 169 (folio 153) with a
% complete sentence ("...is, indeed, a correlation matrix.") and explicitly
% notes no environment is left open. Page 170 (folio 154), the first page
% of this chunk, opens fresh with a subsection heading ("5.2.2 Positive
% definiteness criteria for symmetric matrices"), which confirms nothing
% needs to be closed here -- this file opens directly with \subsection.

\subsection{Positive definiteness criteria for symmetric matrices}

In this section, we summarize the necessary and sufficient conditions established
in the previous sections for a symmetric matrix to be symmetric positive definite
or symmetric positive semidefinite, i.e., the definitions 5.6 and 5.7, the eigenvalues
criteria from Theorem 5.6, Sylvester's Criterion (Theorem 5.8), and we anticipate by
including Cholesky decomposition criteria; cf. Theorem 6.2 from Section 6.1.

Let $A$ be an $n \times n$ symmetric matrix. Equivalent conditions for the matrix $A$
to be symmetric positive definite (spd) or symmetric positive semidefinite (spsd) can
be found in Table 5.1.

\begin{table}[H]
\centering
\caption{Necessary and sufficient conditions for spd and spsd matrices}
\small
\setlength{\tabcolsep}{3pt}
\begin{tabular}{|l c l|c|}
\hline
$A$ spd & $\iff$ & $x\tp Ax > 0, \quad \forall\, x \neq 0$ & definition \\
$A$ spsd & $\iff$ & $x\tp Ax \geq 0, \quad \forall\, x \in \R^n$ & \\
\hline
$A$ spd & $\iff$ & all the eigenvalues of $A$ are $> 0$ & eigenvalue \\
$A$ spsd & $\iff$ & all the eigenvalues of $A$ are $\geq 0$ & criterion \\
\hline
$A$ spd & $\iff$ & all the leading principal minors of $A$ are $> 0$ & Sylvester's \\
$A$ spsd & $\iff$ & all the principal minors of $A$ are $\geq 0$ & criterion \\
\hline
$A$ spd & $\iff$ & the matrix $A$ has Cholesky decomposition &
\begin{tabular}[c]{@{}c@{}}Cholesky\\ decomposition\end{tabular} \\
\hline
\end{tabular}
\end{table}

We include a few comments on the practical relevance of the conditions from
Table 5.1:

\begin{itemize}
\item From a practical standpoint, identifying whether a matrix is symmetric
positive definite is done by running the Cholesky decomposition algorithm from
Table 6.1 on the matrix. If the algorithms breaks down, then the matrix is not
symmetric positive definite; else, the matrix is symmetric positive definite.

\item Computing the eigenvalues of the matrix numerically, e.g., by using the QR
algorithm, in order to find out whether all the eigenvalues of the matrix are
positive, and thus establish whether the matrix is symmetric positive definite, is a
lot more expensive than running the Cholesky algorithm, and potentially imprecise
for matrices with very small positive eigenvalues.

\item Sylvester's Criterion can be applied to establish, without using numerical
methods, whether matrices of small size are symmetric positive definite or symmetric
positive semidefinite.

\item The criteria given by the definitions of spd and spsd matrices are rarely of
practical use; one of the few such examples can be found in Section 5.2.
\end{itemize}

Note that a Cholesky-based criterion also exists for symmetric positive
semidefinite, although it is of little practical importance; see section 6.5 for
details.

\section{The diagonal form of symmetric matrices}

In this section, we include for completeness the proof of the fact that any
symmetric matrix is diagonalizable; see Theorem 5.4. This proof is based on
properties of orthogonal matrices; see section 10.1.3 for details.

\begin{theorem}
Any symmetric matrix is diagonalizable.

More precisely, if $A$ is a symmetric matrix, then there exists an orthogonal matrix
$Q$ and a diagonal matrix $\Lambda$ such that
\begin{equation}
A \;=\; Q\Lambda Q\tp.
\end{equation}
\end{theorem}

\textit{Note that, from Theorem 4.4, it follows that the diagonal entries of the
matrix $\Lambda$ are the eigenvalues of the matrix $A$ and the columns of the matrix
$Q$ are the corresponding eigenvectors of $A$.}

\begin{proof}
We give a proof by induction.

Any $1 \times 1$ matrix is a number, and is the same as a diagonal matrix of size 1.

Assume that any symmetric matrix of size $n - 1$ is diagonalizable. We will show
that any symmetric matrix of size $n$ is also diagonalizable.

Let $A$ be an $n \times n$ symmetric matrix. Let $\lambda_1$ and $v_1$ be an
eigenvalue and a corresponding eigenvector of norm 1 of the matrix $A$, i.e., such
that $Av_1 = \lambda_1v_1$, with $v_1 \neq 0$ and
\begin{equation}
||v_1||^2 \;=\; v_1\tp v_1 \;=\; 1.
\end{equation}

Let $Q_1$ be an $n \times n$ orthogonal matrix with the vector $v_1$ as the first
column, i.e., let
\[
Q_1 \;=\; (v_1 \mid q_2 \mid \dots \mid q_n)
\]
where $q_i$ are $n \times 1$ column vectors such that
\begin{equation}
(v_1,q_i) \;=\; q_i\tp v_1 \;=\; 0, \quad \forall\, i = 2:n,
\end{equation}
and $||q_i||^2 = q_i\tp q_i = 1$, for $i = 2:n$.

From the matrix-matrix multiplication formula (1.11), and using the fact that
$Av_1 = \lambda_1v_1$, we find that
\[
AQ_1 \;=\; (Av_1 \mid Aq_2 \mid \dots \mid Aq_n) \;=\; (\lambda_1v_1 \mid Aq_2 \mid
\dots \mid Aq_n).
\]

Then,
\begin{equation}
\begin{split}
Q_1\tp AQ_1 \;&=\;
\begin{pmatrix} v_1\tp\\ --\\ q_2\tp\\ --\\ \vdots\\ --\\ q_n\tp \end{pmatrix}
(\lambda_1v_1 \mid Aq_2 \mid \dots \mid Aq_n)\\
&=\;
\begin{pmatrix}
\lambda_1v_1\tp v_1 & v_1\tp Aq_2 & \dots & v_1\tp Aq_n\\
\lambda_1q_2\tp v_1 & q_2\tp Aq_2 & \dots & q_2\tp Aq_n\\
\vdots & \vdots & & \vdots\\
\lambda_1q_n\tp v_1 & q_n\tp Aq_2 & \dots & q_n\tp Aq_n
\end{pmatrix}\\
&=\;
\begin{pmatrix}
\lambda_1 & v_1\tp Aq_2 & \dots & v_1\tp Aq_n\\
0 & q_2\tp Aq_2 & \dots & q_2\tp Aq_n\\
\vdots & \vdots & & \vdots\\
0 & q_n\tp Aq_2 & \dots & q_n\tp Aq_n
\end{pmatrix},
\end{split}
\end{equation}
where (5.66) and (5.67) were used for the last equality.

Since the matrix $A$ is symmetric, $A\tp = A$, and the matrix $Q_1\tp AQ_1$ is also
symmetric:
\begin{equation}
(Q_1\tp AQ_1)\tp \;=\; Q_1\tp A\tp (Q_1\tp)\tp \;=\; Q_1\tp AQ_1.
\end{equation}
Then, from (5.68), it follows that all the entries $(1,k)$, with $2 \leq k \leq n$,
from the first row of the matrix $Q_1\tp AQ_1$ must also be equal to 0, i.e.,
\[
Q_1\tp AQ_1 \;=\;
\begin{pmatrix}
\lambda_1 & 0 & \dots & 0\\
0 & q_2\tp Aq_2 & \dots & q_2\tp Aq_n\\
\vdots & \vdots & & \vdots\\
0 & q_n\tp Aq_2 & \dots & q_n\tp Aq_n
\end{pmatrix}.
\]

Denote by $A_{n-1}$ the $(n-1) \times (n-1)$ matrix made of the rows $2:n$ and the
columns $2:n$ of the matrix $Q_1\tp AQ_1$, i.e.,
\[
A_{n-1} \;=\;
\begin{pmatrix}
q_2\tp Aq_2 & \dots & q_2\tp Aq_n\\
\vdots & \vdots & \vdots\\
q_n\tp Aq_2 & \dots & q_n\tp Aq_n
\end{pmatrix}.
\]

Then,
\begin{equation}
Q_1\tp AQ_1 \;=\;
\begin{pmatrix}
\lambda_1 & 0 & \dots & 0\\
0 & & & \\
\vdots & & A_{n-1} & \\
0 & & &
\end{pmatrix}.
\end{equation}

Since $Q_1\tp AQ_1$ is symmetric, it follows that $A_{n-1}$ is a symmetric matrix of
size $n - 1$. Then, from the induction hypothesis, we obtain that $A_{n-1}$ is
diagonalizable, i.e., there exist an orthogonal matrix $Q_{n-1}$ of size $n-1$ and a
diagonal matrix $\Lambda_{n-1}$ of size $n-1$ such that
\begin{equation}
A_{n-1} \;=\; Q_{n-1}\Lambda_{n-1}Q_{n-1}\tp.
\end{equation}
From (5.70) and (5.71), we obtain that
\[
\begin{aligned}
Q_1\tp AQ_1 \;&=\;
\begin{pmatrix}
\lambda_1 & 0 & \dots & 0\\
0 & & & \\
\vdots & & A_{n-1} & \\
0 & & &
\end{pmatrix}
\;=\;
\begin{pmatrix}
\lambda_1 & 0 & \dots & 0\\
0 & & & \\
\vdots & & Q_{n-1}\Lambda_{n-1}Q_{n-1}\tp & \\
0 & & &
\end{pmatrix}\\[6pt]
&=\;
\begin{pmatrix}
1 & 0 & \dots & 0\\
0 & & & \\
\vdots & & Q_{n-1} & \\
0 & & &
\end{pmatrix}
\begin{pmatrix}
\lambda_1 & 0 & \dots & 0\\
0 & & & \\
\vdots & & \Lambda_{n-1} & \\
0 & & &
\end{pmatrix}
\begin{pmatrix}
1 & 0 & \dots & 0\\
0 & & & \\
\vdots & & Q_{n-1}\tp & \\
0 & & &
\end{pmatrix},
\end{aligned}
\]
which can be written as
\begin{equation}
Q_1\tp AQ_1 \;=\; Q_n\Lambda_nQ_n\tp,
\end{equation}
where $Q_n = \begin{pmatrix} 1 & 0 & \dots & 0\\ 0 & & & \\ \vdots & & Q_{n-1} & \\
0 & & & \end{pmatrix}$ and $\Lambda_n = \begin{pmatrix} \lambda_1 & 0 & \dots & 0\\
0 & & & \\ \vdots & & \Lambda_{n-1} & \\ 0 & & & \end{pmatrix}$ are an orthogonal
matrix of size $n$ and a diagonal matrix of size $n$, respectively.

We multiply (5.72) to the left by the matrix $Q_1$ and to the right by the matrix
$Q_1\tp$. Since $Q_1Q_1\tp = I$, see (10.14), it follows that
\[
\begin{aligned}
Q_1(Q_1\tp AQ_1)Q_1\tp &= Q_1(Q_n\Lambda_nQ_n\tp)Q_1\tp\\
\iff (Q_1Q_1\tp)A(Q_1Q_1\tp) &= (Q_1Q_n)\Lambda_n(Q_1Q_n)\tp\\
\iff A &= Q\Lambda_nQ\tp,
\end{aligned}
\]
where $Q = Q_1Q_n$.

Note that $Q$ is an orthogonal matrix, since it is the product of two orthogonal
matrices; cf. Lemma 10.6. Then, $Q\tp = Q^{-1}$, see (10.15), and
\[
A \;=\; Q\Lambda_nQ\tp \;=\; Q\Lambda_nQ^{-1}
\]
is the diagonal form of $A$. Thus, the matrix $A$ is diagonalizable.

We conclude, by induction, that any symmetric matrix is diagonalizable.
\end{proof}

\section{References}

An elegant proof for Sylvester's criterion can be found in Gilbert \cite{16}.

% PDF page 173 (folio 157) ends the References section here; the rest of
% the page is blank. No environment is left open. Page 174 (folio 158)
% opens fresh with "5.5 Exercises", continued in sec05_c3_2.tex.
